\documentclass{article}

% -----------------------------
% Packages
% -----------------------------
\usepackage{fancyhdr}
\usepackage{extramarks}
\usepackage{amsmath}
\usepackage{amsthm}
\usepackage{amsfonts}
\usepackage{tikz}
\usepackage[plain]{algorithm}
\usepackage{algpseudocode}
\usepackage{amssymb}
\usepackage{enumitem}

\usetikzlibrary{automata,positioning}

% -----------------------------
% Basic Document Settings
% -----------------------------
\topmargin=-0.45in
\evensidemargin=0in
\oddsidemargin=0in
\textwidth=6.5in
\textheight=9.0in
\headsep=0.25in

\linespread{1.1}

\pagestyle{fancy}
\lhead{\hmwkAuthorName}
\chead{\hmwkClass: \hmwkTitle}
\rhead{}
\lfoot{\lastxmark}
\cfoot{\thepage}

\renewcommand\headrulewidth{0.4pt}
\renewcommand\footrulewidth{0.4pt}

\setlength\parindent{0pt}

% -----------------------------
% Problem Section Helpers
% -----------------------------
\newcommand{\enterProblemHeader}[1]{
  \nobreak\extramarks{}{Problem #1 continued on next page\ldots}\nobreak{}
  \nobreak\extramarks{Problem #1 (continued)}{Problem #1 continued on next page\ldots}\nobreak{}
}

\newcommand{\exitProblemHeader}[1]{
  \nobreak\extramarks{Problem #1 (continued)}{Problem #1 continued on next page\ldots}\nobreak{}
  \nobreak\extramarks{Problem #1}{}\nobreak{}
}

\setcounter{secnumdepth}{0}
\newcounter{partCounter}
\newcounter{homeworkProblemCounter}
\setcounter{homeworkProblemCounter}{1}
\nobreak\extramarks{Problem \arabic{homeworkProblemCounter}}{}\nobreak{}

% -----------------------------
% Theorem env (optional)
% -----------------------------
\newtheorem*{theorem}{Theorem}

% -----------------------------
% Homework Problem Environment
%   Usage:
%   \begin{homeworkProblem}            % auto-number
%     % Your answer here
%   \end{homeworkProblem}
%
%   \begin{homeworkProblem}[1.2.3]     % label problem "1.2.3" (e.g. textbook numbering)
%     % Your answer here
%   \end{homeworkProblem}
% -----------------------------
\newenvironment{homeworkProblem}[1][]{
  \def\hwProblemLabel{#1}
  \ifx\hwProblemLabel\empty
    \edef\hwProblemLabel{\arabic{homeworkProblemCounter}}
    \stepcounter{homeworkProblemCounter}
  \fi
  \section{Problem \hwProblemLabel}
  \setcounter{partCounter}{1}
  \enterProblemHeader{\hwProblemLabel}
}{
  \exitProblemHeader{\hwProblemLabel}
}

% -----------------------------
% Homework Details (EDIT THESE)
% -----------------------------
\newcommand{\hmwkTitle}{Homework \#6}
\newcommand{\hmwkClass}{Math 3320}
\newcommand{\hmwkAuthorName}{\textbf{Joseph Quinn}}

% -----------------------------
% Title
% -----------------------------
\title{
  \vspace{2in}
  \textmd{\textbf{\hmwkClass:\ \hmwkTitle}}\\
  \vspace{3in}
}
\author{\hmwkAuthorName}
\date{}

% -----------------------------
% Convenience Macros (optional)
% -----------------------------
\renewcommand{\part}[1]{\textbf{\large Part \Alph{partCounter}}\stepcounter{partCounter}\\}

% Algorithms
\newcommand{\alg}[1]{\textsc{\bfseries \footnotesize #1}}

% Calculus
\newcommand{\deriv}[1]{\frac{\mathrm{d}}{\mathrm{d}x} (#1)}
\newcommand{\pderiv}[2]{\frac{\partial}{\partial #1} (#2)}
\newcommand{\dx}{\mathrm{d}x}

% Statistics
\newcommand{\E}{\mathrm{E}}
\newcommand{\Var}{\mathrm{Var}}
\newcommand{\Cov}{\mathrm{Cov}}
\newcommand{\Bias}{\mathrm{Bias}}

% Proof step (for aligned derivations)
\newcommand{\step}[2]{& #1 & & \text{#2} \\}

% Solution header (optional)
\newcommand{\solution}{\textbf{\large Solution}}

% -----------------------------
% Document
% -----------------------------
\begin{document}

\maketitle
\pagebreak
\begin{homeworkProblem}[3.1.6]

	\begin{enumerate}
		\item[(a)]
		      \( n=15 \) and \( k=3 \)

		      to find \( d \) we find the min weight nonzero codeword which is 5 so \( d=5 \)



		      Thus, this code can correct all errros of \( \le t \) with,
		      \[
			      t=\left\lfloor\frac{d-1}{2}\right\rfloor
			      =\left\lfloor\frac{5-1}{2}\right\rfloor=2.
		      \]

		      (\( \mid C \mid = 2^k \) so we can sub that in for it)

		      So we can now substitute into the hamming bound formula

		      \[
			      \begin{aligned}
				      2^k\sum_{i=0}^{t}\binom{n}{i}  & \le 2^n    \\
				      2^3\sum_{i=0}^{2}\binom{15}{i} & \le 2^{15} \\
				      8\left(1+15+105\right)         & \le 32768  \\
				      968\le 32768
			      \end{aligned}
		      \]
		      This checks out so its valid bound

		\item[(b)]
		      Here $n=6$, $k=3$, and $d=3$, making our \( t \) the following,
		      \[
			      t=\left\lfloor\frac{3-1}{2}\right\rfloor=1.
		      \]

		      By the Hamming bound,
		      \[
			      \begin{aligned}
				      2^k\sum_{i=0}^{t}\binom{n}{i} & \le 2^n \\
				      2^3\sum_{i=0}^{1}\binom{6}{i} & \le 2^6 \\
				      8(1+6)                        & \le 64  \\
				      56\le 64
			      \end{aligned}
		      \]

		      So, the Hamming bound is satisfied.

		\item[(c)]

		      Here $n=7$, $k=4$, and $d=3$. with \( t \) being calculated as the following,
		      \[
			      t=\left\lfloor\frac{3-1}{2}\right\rfloor=1.
		      \]

		      By the Hamming bound,
		      \[
			      \begin{aligned}
				      2^k\sum_{i=0}^{t}\binom{n}{i} & \le 2^n \\
				      2^4\sum_{i=0}^{1}\binom{7}{i} & \le 2^7 \\
				      16(1+7)                       & \le 128 \\
				      128\le 128
			      \end{aligned}
		      \]

		      So, the Hamming bound is satisfied (with bonus of being a perfect code!)

	\end{enumerate}
\end{homeworkProblem}

\begin{homeworkProblem}[3.1.11]
	Because \( k \) columns are linearly dependent this means a linear combination of them equals 0, since \( c=vG \) creates a codeword through the generator matrix \(G \),
	this is the same as taking a linear combination of the columns in \( G \) that align with \( 1 \)'s in \( v \). So this tells us there is some nonzero word that results in \( 0 \)'s from the depdentet columhns. If we use this same word
	in \( c=vG \) (assuming the columns are in hte same place) we will get a codeword \( c \) where the depdent columns are set to 0.
\end{homeworkProblem}

\begin{homeworkProblem}[3.1.21]

	We want a $(15,6,5)$ code, so the parity-check matrix
	$H$ must have $15$ rows and $15-6=9$ columns.
	\\
	We know thatevery set of $4$ rows
	must be linearly independent, and some set of
	$5$ rows must be dependent.

	From the example we can start with the identity matrix as the first 9 rows
	of $H$.
	\\
	The example then gives us another 3 rows to work with,
	\[
		\begin{bmatrix}
			I_9       \\
			111100000 \\
			100011100 \\
			101000011
		\end{bmatrix}
	\]
	So we only need to find the last $3$ rows.

	\\

	Suppose $j$ of the added (non-identity) rows sum to a
	vector of weight $m$. That vector is also a sum of $m$ rows of
	$I_9$, so together that is a dependence among $j+m$ rows.
	For no $4$ rows to be dependent we need $j+m\ge 5$, i.e.
	any sum of $j$ added rows must have weight at least $5-j$:
	This means any new row we had must have a weight \( \ge 4 \), the sum of any $2$ added rows has weight $\ge 3$,
	the sum of any $3$ added rows has weight $\ge 2$, and the sum of any $4$ added rows is nonzero.
	\\
	The rows $010101010$, $000001111$, $000110011$ satisfy these
	conditions together with the example's rows, giving us a valid \( H \) matrix.
	\[
		H=
		\begin{bmatrix}
			I_9       \\
			111100000 \\
			100011100 \\
			101000011 \\
			010101010 \\
			000001111 \\
			000110011
		\end{bmatrix}.
	\]

	So $H$ defines a $(15,6,5)$ code.

\end{homeworkProblem}

\begin{homeworkProblem}[3.2.5]
	\begin{proof}
		We have the following sum of binomials
		\[
			\binom{n}{0}+\binom{n}{1}
		\]

		\( n \) choose 0 of any number \( n \) is always \( 1 \) as there is only one way to choose nothing, and \( n \) choose 1 is always \( n \)
		as there is always \( n \) ways to choose a single item from a set of \( n \).

		So

		\[
			\binom{n}{0}+\binom{n}{1} = 1 + n
		\]
		Then if we sub in $n=2^r-1$, we have

		\[
			\begin{aligned}
				 & =1+n       \\
				 & =1+(2^r-1) \\
				 & =2^r.
			\end{aligned}
		\]
	\end{proof}

\end{homeworkProblem}
\begin{homeworkProblem}
	\begin{enumerate}
		\item[(a)]

		      $n=15$, $d=3$

		      Since $d=3$, we have $t=1$. Thus,
		      \[
			      \begin{aligned}
				      |C|\left(\binom{15}{0}+\binom{15}{1}\right)
				                & =2^{15}                    \\
				      |C|(1+15) & =2^{15}                    \\
				      |C|       & =\frac{2^{15}}{16}=2^{11}.
			      \end{aligned}
		      \]

		      If we look at the hamming code with \( r= 4 \) we get:
		      \begin{itemize}
			      \item \( n=2^4-1=15 \)
			      \item \( d=3 \)
			      \item \( k=2^4-1-4=16-1-5=11\)
		      \end{itemize}

		      Because \( k=11 \), \( \mid C \mid = 2^11 \) which is exactly our hamming bound!

		      Meaning this hamming code is a perfect code satisfying the above requirements!


		\item[(b)]
		      $n=31$, $d=3$

		      Again, $t=1$. Thus,
		      \[
			      \begin{aligned}
				      |C|\left(\binom{31}{0}+\binom{31}{1}\right)
				                & =2^{31}                    \\
				      |C|(1+31) & =2^{31}                    \\
				      |C|       & =\frac{2^{31}}{32}=2^{26}.
			      \end{aligned}
		      \]

		      If we look at the hamming code with \( r= 5\) we get:
		      \begin{itemize}
			      \item \( n=2^5-1 = 31 \)
			      \item \( d=3 \)
			      \item \( k=2^5-1-5=32-1-5=26\)
		      \end{itemize}

		      Because \( k=26 \), \( \mid C \mid = 2^26 \) which is exactly our hamming bound!

		      Meaning this hamming code is a perfect code satisfying the above requirements!

		\item[(c)]
		      $n=15$, $d=5$

		      Since $d=5$, we have $t=2$. Thus,
		      \[
			      \begin{aligned}
				      |C|\left(
				      \binom{15}{0}+
				      \binom{15}{1}+
				      \binom{15}{2}
				      \right)       & =2^{15}              \\
				      |C|(1+15+105) & =2^{15}              \\
				      |C|           & =\frac{2^{15}}{121}.
			      \end{aligned}
		      \]

		      Since the denominator is not a power of two the fraction is not an integer meaning no code will satisfy this condition.

	\end{enumerate}
\end{homeworkProblem}


\begin{homeworkProblem}[3.3.3]
	For a binary Hamming code of length $n=15$, we have
	\(
	n=2^r-1.
	\)
	meaning
	\[
		15=2^4-1 \implies r=4.
	\]

	With this, we can find the dimension:
	\[
		k=n-r=15-4=11.
	\]

	Our generator matrix must be $11\times 15$, and
	in standard form it will look like,
	\[
		G=[I_{11}\mid P].
	\]

	First to get here we must construct the parity check matrix as we need to maintain that each row is non zero.

	\[
		H=
		\begin{bmatrix}
			P \\
			I_4
		\end{bmatrix}.
	\]

	Its fairly trivial to fill the matrix as there are $2^4=16$ binary vectors of length 4 (15 minus the zero vector).
	so we can just fill in all the vectors as rows making sure to keep the identity matrix at the bottom.

	\[
		H=
		\left[
			\begin{array}{cccc}
				0 & 0 & 1 & 1 \\
				0 & 1 & 0 & 1 \\
				0 & 1 & 1 & 0 \\
				0 & 1 & 1 & 1 \\
				1 & 0 & 0 & 1 \\
				1 & 0 & 1 & 0 \\
				1 & 0 & 1 & 1 \\
				1 & 1 & 0 & 0 \\
				1 & 1 & 0 & 1 \\
				1 & 1 & 1 & 0 \\
				1 & 1 & 1 & 1 \\ \hline
				1 & 0 & 0 & 0 \\
				0 & 1 & 0 & 0 \\
				0 & 0 & 1 & 0 \\
				0 & 0 & 0 & 1
			\end{array}
			\right]
		=
		\begin{bmatrix}
			P \\
			I_4
		\end{bmatrix}.
	\]

	This is in standard form so to get the associated generator matrix we simply keep \( P \) as is and add the \( I_{11} \) identity matrix to the left.
	we obtain
	\[
		G=
		\left[
			\begin{array}{c|cccc}
				       & 0 & 0 & 1 & 1 \\
				       & 0 & 1 & 0 & 1 \\
				       & 0 & 1 & 1 & 0 \\
				       & 0 & 1 & 1 & 1 \\
				I_{11} & 1 & 0 & 0 & 1 \\
				       & 1 & 0 & 1 & 0 \\
				       & 1 & 0 & 1 & 1 \\
				       & 1 & 1 & 0 & 0 \\
				       & 1 & 1 & 0 & 1 \\
				       & 1 & 1 & 1 & 0 \\
				       & 1 & 1 & 1 & 1
			\end{array}
			\right].
	\]

	Now to encode the message,
	\(
	u=11111000000.
	\)
	with
	\[
		c=uG.
	\]

	Since its in standard form the first 11 digits will simply be \( u \), and the remaining 4 will be the linear combinations of the last 4 columns, since the first
	5 digits of \( u \) are 1 its easy and is simply the sum of the first 5 rows.

	\[
		\begin{aligned}
			 & =0011+0101+0110+0111+1001 \\
			 & =1110.
		\end{aligned}
	\]

	and finally add both pieces together,
	\[
		c=[11111000000\mid 1110].
	\]

	making the final encoded message,
	\[
		c=111110000001110.
	\]
\end{homeworkProblem}

\begin{homeworkProblem}[3.3.5]
	We use the same parity check matrix \( H \) from the last problem

	Since \( d=3 \), the code corrects any single error, so the coset leaders are \( 0 \) and the
	15 words of weight 1.
	\[
		\begin{array}{c|c}
			\text{Syndrome} & \text{Coset Leader} \\
			\hline
			0000            & 000000000000000     \\
			0011            & 100000000000000     \\
			0101            & 010000000000000     \\
			0110            & 001000000000000     \\
			0111            & 000100000000000     \\
			1001            & 000010000000000     \\
			1010            & 000001000000000     \\
			1011            & 000000100000000     \\
			1100            & 000000010000000     \\
			1101            & 000000001000000     \\
			1110            & 000000000100000     \\
			1111            & 000000000010000     \\
			1000            & 000000000001000     \\
			0100            & 000000000000100     \\
			0010            & 000000000000010     \\
			0001            & 000000000000001
		\end{array}
	\]

	To decode \( w \) we compute \( wH \) (the sum of the rows of \( H \) where \( w \) has a 1),
	look up its coset leader \( u \), and decode to \( c=w+u \).

	\begin{enumerate}
		\item[(a)]
		      \( w=01010\ 01010\ 01000 \) has 1's in positions \( 2,4,7,9,12 \), so
		      \[
			      \begin{aligned}
				      wH & =0101+0111+1011+1101+1000 \\
				         & =1100.
			      \end{aligned}
		      \]
		      From the SDA, \( 1100 \) has coset leader \( u=000000010000000 \), so
		      \[
			      c=w+u=01010\ 01110\ 01000.
		      \]

		\item[(b)]
		      \( w=11110\ 00101\ 10110 \) has 1's in positions \( 1,2,3,4,8,10,11,13,14 \), so
		      \[
			      \begin{aligned}
				      wH & =0011+0101+0110+0111+1100+1110+1111+0100+0010 \\
				         & =1100.
			      \end{aligned}
		      \]
		      Again the coset leader is \( u=000000010000000 \), so
		      \[
			      c=w+u=11110\ 00001\ 10110.
		      \]

		\item[(c)]
		      \( w=11100\ 01110\ 00111 \) has 1's in positions \( 1,2,3,7,8,9,13,14,15 \), so
		      \[
			      \begin{aligned}
				      wH & =0011+0101+0110+1011+1100+1101+0100+0010+0001 \\
				         & =1101.
			      \end{aligned}
		      \]
		      From the SDA, \( 1101 \) has coset leader \( u=000000001000000 \), so
		      \[
			      c=w+u=11100\ 01100\ 00111.
		      \]
	\end{enumerate}
\end{homeworkProblem}
\end{document}
